\begin{figure}[t]
{
    \centering
    \includegraphics[width=\textwidth]{figures/motivation_f.pdf}
    \caption{
    \textbf{Overview of the key differences between I-JEPA and \ours.} I-JEPA predicts embeddings of randomly sampled target regions in a parallel and independent manner. In contrast, \ours~ranks regions by attention and predicts region embeddings sequentially from high- to low-attention regions. Patches shown are only for illustration; the model predicts region embeddings rather than raw pixels.
    }
    \label{fig:motivation}
}
\end{figure}



% In text
% \begin{wrapfigure}{r}{0.57\textwidth}
%     \centering
%     % \vspace{-1.7\baselineskip} % optional: tighten top spacing
%     \includegraphics[width=0.57\textwidth]{figures/human_attention_inspired.pdf}
%     \setlength{\abovecaptionskip}{0.1cm}
%     \caption{\textbf{(A) Humans} perceive visual scenes selectively and sequentially, focusing on discriminative regions.
%     % such as the red head and red chest. 
%     \textbf{(B) \ours} emulates this process by ranking regions by attention-derived importance and successively predicting next discriminative region embedding.}
%     \label{fig:human}
%     \vspace{-1.8\baselineskip} % optional: tighten bottom spacing
% \end{wrapfigure}